\chapter{Gas Safety and Hot Work Risk Assessment}
\label{chap:Gas Safety and Hot Work Risk Assessment}

%---------------------------- SECTION ----------------------------
\section{Why this assessment matters in construction}

Gas incidents on construction sites fall into two categories: inadvertent contact with existing gas infrastructure during excavation or structural alteration, and planned hot work activities (welding, cutting, brazing, and use of LPG torches) that create ignition sources in proximity to flammable or combustible materials. Both categories have produced multiple fatalities on UK construction sites in recent years and both are managed by regulatory frameworks that are well established but inconsistently applied.

The Gas Safety (Installation and Use) Regulations 1998 (GSIUR) and the Gas Act 1986 establish the framework for work on gas systems in the UK. All gas work on domestic, commercial, and industrial gas systems must be carried out by a Gas Safe registered engineer; working on a gas system without Gas Safe registration is a criminal offence and a significant enforcement priority for the Health and Safety Executive and the Gas Safe Register. Despite the clarity and well-publicised nature of this requirement, unlicensed gas work by unregistered persons --- ranging from the illegal connection of a gas appliance by a general plumber to the modification of a gas main by a groundworks contractor who did not recognise the pipe as gas --- continues to account for a significant proportion of gas incidents reported under RIDDOR.

Hot work --- any process that generates heat, sparks, or open flame in proximity to combustible materials --- is the cause of a significant proportion of construction site fires. Welding, disc cutting, angle grinding, LPG torch application, and oxy-fuel cutting all generate sparks or radiant heat that can ignite adjacent combustible materials including dry timber, insulation, packaging, and roof felt. The consequences of a construction site fire can include: loss of the building under construction; extension of fire to adjacent properties; injury or death of operatives who cannot escape; and the disproportionate impact on the programme and finances of the project. A hot work permit system is the administrative control mechanism through which the ignition risk is managed; its effectiveness depends on the quality of the pre-work inspection, the clearance of combustibles, the provision of fire watch, and the one-hour post-work fire watch period after hot work has ceased.

Construction sites using LPG (liquefied petroleum gas) for heating, hot work, and equipment operation introduce an additional hazard: LPG is heavier than air and will accumulate in low-lying areas, basements, and drainage runs where it can form a flammable mixture that ignites from the most trivial spark source. LPG cylinder storage, handling, and connection requirements are specific and must be addressed in the risk assessment.

\barquote{``A hot work permit that is filled in after the work has started is a paperwork exercise, not a risk control. The value of the permit is in the pre-work inspection and clearance that it mandates.''}

\example{Site Reality}{A roofing contractor is waterproofing a flat commercial roof using a bitumen torch-on felt system. The flat roof has a 150\,mm raised kerb at the junction with the adjoining pitched roof. Below the kerb, within the roof build-up, is a layer of expanded polystyrene insulation. The roofer applies the LPG torch to the felt at the kerb junction; the heat travels through the kerb detail into the polystyrene, which ignites below the visible surface. The roofer completes the kerb detail and leaves the roof. There is no post-work fire watch. Forty minutes later, smoke is visible from the kerb junction; the polystyrene fire has spread laterally beneath the roof build-up. The roof structure, a section of the main roof deck, and the accommodation below suffer significant fire damage. The hot work risk assessment specified a ``30-minute fire watch after work''; the fire watch was not conducted.}

%---------------------------- SECTION ----------------------------
\section{Legal and professional framework}

The Gas Safety (Installation and Use) Regulations 1998 are the primary legislative instrument for gas work safety, requiring that all gas work is carried out by a registered Gas Safe engineer. The Dangerous Substances and Explosive Atmospheres Regulations 2002 (DSEAR) apply where flammable gases (including LPG and natural gas) are stored or handled on construction sites. The Highly Flammable Liquids and Liquefied Petroleum Gas Regulations 1972 (as applied) provide specific requirements for LPG storage. The Health and Safety at Work etc.\ Act 1974 and MHSWR 1999 impose the general duty framework; CDM 2015 requires the Construction Phase Plan to address hot work management.

\paragraph{Key frameworks relevant to this chapter include: GSIUR 1998; DSEAR 2002; MHSWR 1999; CDM 2015; LPG storage regulations.}

\begin{itemize}
  \bitem{Gas Safety (Installation and Use) Regulations 1998}{Require that all work on gas fittings (installation, repair, maintenance, disconnection) is carried out by a Gas Safe registered engineer; make it a criminal offence to carry out gas work without Gas Safe registration; apply to all gas work on domestic, commercial, and industrial gas systems.}
  \bitem{DSEAR 2002}{Apply wherever a dangerous substance (including LPG, acetylene, and natural gas) may be present on a construction site; require that explosive atmosphere zones are classified, that ignition sources are controlled, that emergency procedures are in place, and that workers are trained in the risks and controls.}
  \bitem{CDM 2015 Construction Phase Plan}{Must address the management of hot work on the project, including the hot work permit system, the fire watch requirements, the storage and use of LPG, the management of confined spaces with gas risk, and the emergency procedures for gas escape or fire.}
  \bitem{Fire Safety Order 2005 (as applicable)}{The Regulatory Reform (Fire Safety) Order 2005 applies to the non-domestic parts of construction sites used as workplaces; fire risk associated with hot work, LPG, and combustible materials must be reflected in the fire risk assessment for the site.}
\end{itemize}

%---------------------------- SECTION ----------------------------
\section{Conducting the assessment using the five-step method}

\subsection{Step 1 --- Identify hazards}
\paragraph{Gas and hot work hazards in construction include: unintended contact with live gas mains during excavation causing explosion and fire; escape of gas from damaged or improperly connected gas fittings during plumbing and mechanical installation works; fire from hot work sparks or radiant heat igniting adjacent combustible materials; subsurface fire in roof or floor build-ups ignited by torch-on application through the building envelope; LPG accumulation in low-lying areas, basements, or plant rooms from poorly stored or connected cylinders; explosion from LPG accumulation in an enclosed space; acetylene cylinder fire from reverse flame entry during oxy-fuel cutting; fire from oxygenated spaces adjacent to welding operations.}

\subsection{Step 2 --- Decide who or what is affected}
\paragraph{Persons at risk include: all operatives on or in the building during hot work activities; gas engineers conducting connection or disconnection work on live gas systems; operatives in areas where LPG may accumulate; operatives in adjacent areas who may be affected by a fire initiated by hot work; members of the public adjacent to building work where LPG is stored or used; emergency services attending a gas escape or fire incident; and occupants of adjacent buildings affected by a fire originating on the construction site.}

\subsection{Step 3 --- Evaluate likelihood and consequence}
\paragraph{The consequence of an LPG explosion or fire on a construction site is Very High (fatality and major structural damage). The likelihood of a hot work fire where no permit system or fire watch is in place is Medium to High, given the frequency of hot work on construction sites and the prevalence of combustible building materials in the vicinity of the work. Inherent risk for uncontrolled hot work must be rated High to Very High (16--25/25). A properly implemented hot work permit with pre-work combustible clearance, fire suppression provision, and one-hour post-work fire watch reduces residual risk to Low (4--6/25).}

\subsection{Step 4 --- Select and implement controls}
\paragraph{Hot work controls: implement a formal hot work permit system with pre-work inspection by the site manager or supervisor; clear combustibles to a minimum of 10 metres from the hot work location; cover unavoidable combustibles with fire-retardant welding blankets; provide a dry powder or CO$_2$ extinguisher at the hot work location; maintain a trained fire watch during all hot work and for minimum one hour after all hot work has ceased; check the hot work area and all adjacent areas (including above and below the work level) at the end of the fire watch period. Gas work controls: confirm Gas Safe registration of all gas operatives; use a permit-to-work for any work on or near live gas systems; test gas lines before reconnection; install and test gas leak detector in enclosed spaces where gas is used.}

\subsection{Step 5 --- Record and review}
\paragraph{Hot work permit records must be retained for the duration of the project. Gas Safe registration certificates for all gas operatives must be checked before work commences and retained on file. LPG cylinder delivery and storage records must be maintained. Any gas escape, hot work fire, or fire watch finding must be recorded and reviewed.}

\example{Five-Step Application}{A mechanical and electrical contractor is installing a new gas-fired boiler plant in a commercial building. Step 1 identifies: gas connection to live gas main by mechanical operative; LPG torch use for copper pipe soldering adjacent to timber roof structure; LPG cylinder storage on scaffolding. Step 2 identifies: mechanical operative; gas safe registered engineer; roofer working above the boiler room. Step 3 rates inherent gas connection risk as Very High (25/25) if performed without Gas Safe registration. Step 4 implements: Gas Safe registered engineer performing all gas connections; hot work permit issued for all LPG torch work; combustibles cleared from 5-metre radius of torch work; welding blanket used to protect timber roof structure at eaves; fire watch maintained for one hour post-torch work; LPG cylinders stored in external ventilated cage, not on scaffold. Residual risk rated Low (6/25). Step 5: Gas Safe registration confirmed before site commencement.}

%---------------------------- SECTION ----------------------------
\section{Applying the hierarchy of controls}

\paragraph{The hierarchy for gas and hot work risk places avoidance and substitution of hot work at the highest level, supported by engineering controls (fire suppression, gas detection) and a comprehensive permit system.}

\begin{itemize}
  \bitem{Elimination}{Avoid hot work entirely where design alternatives exist: specify push-fit plumbing systems that eliminate soldering; specify cold-applied bitumen systems for flat roofs rather than torch-on; specify prefabricated pipework assemblies delivered to site for mechanical connection rather than site welding.}
  \bitem{Substitution}{Replace oxy-fuel cutting with disc cutting or plasma cutting where structurally appropriate; replace LPG torch soldering with press-fit or push-fit pipe connections; replace bituminous torch-on waterproofing with cold-applied systems.}
  \bitem{Engineering controls}{Fire suppression system (extinguisher) at the hot work location; fire detection within the area of hot work; gas detection in enclosed spaces where gas work is conducted; welding blankets and fire-retardant barriers around unavoidable combustibles; ventilation of enclosed spaces before LPG use.}
  \bitem{Administrative controls}{Hot work permit system with pre-work inspection and combustible clearance; one-hour post-work fire watch; Gas Safe registration verification; DSEAR zone classification for LPG storage; emergency gas isolation procedure; LPG cylinder inspection and storage compliance; permit-to-work for work on live gas systems.}
  \bitem{PPE}{Welding visor and flash-safe goggles for welding and cutting; fire-retardant overalls for torch and welding operatives; welding gloves; safety boots; hearing protection for grinding and disc cutting operations.}
\end{itemize}

%---------------------------- SECTION ----------------------------
\section{Cadence of assessment and review triggers}

\paragraph{The gas and hot work risk assessment must be reviewed as the types and locations of hot work change through the project, and the hot work permit system provides the ongoing assessment mechanism for each individual operation.}

\begin{itemize}
  \bitem{Before each hot work permit is issued}{The site manager or supervisor must conduct a pre-work inspection of the proposed hot work location before each permit is issued; this is the operational mechanism of the risk assessment, not a separate activity.}
  \bitem{Before any gas operative commences work}{Gas Safe registration must be confirmed for each gas operative before they commence any gas work on site; the registration must be for the category of gas work they are undertaking.}
  \bitem{Following any hot work fire or near-miss}{Any fire, smouldering, or smoke incident associated with hot work must trigger an immediate review of the permit system and the controls applied; work must not resume until the review is complete.}
  \bitem{When new LPG equipment is brought to site}{Any new LPG cylinder, appliance, or distribution arrangement must be reviewed against the DSEAR zone classification and the LPG storage requirements before use.}
  \bitem{At minimum annually}{The overall gas and hot work risk assessment for ongoing projects must be reviewed annually to confirm currency with current activities, materials, and regulatory requirements.}
\end{itemize}

%---------------------------- SECTION ----------------------------
\section{Common failure modes}

\begin{itemize}
  \bitem{No fire watch conducted after hot work}{The hot work permit specifies a post-work fire watch period; the fire watch is not conducted because work pressure prevents the operative from remaining at the location after the visible hot work has ceased. Sub-surface fires in roof and wall build-ups are the predictable consequence.}
  \bitem{Gas Safe registration not verified}{Gas work is conducted by general plumbers or heating engineers who hold Gas Safe registration for gas appliance work but not for the specific category of work they are performing; the principal contractor does not verify the category of registration against the work to be done.}
  \bitem{Combustibles not cleared before hot work}{Hot work is conducted with timber, insulation, or packaging within the combustibles clearance zone specified in the permit; the pre-work inspection does not adequately identify the combustibles present.}
  \bitem{LPG stored incorrectly}{LPG cylinders stored in basement areas, enclosed stores, or near to ignition sources; cylinders not in use stored without valve protection caps; connection hoses not inspected or replaced at prescribed intervals.}
  \bitem{Permit system bypassed}{Operatives conduct hot work without obtaining a permit because the permit system is perceived as bureaucratic and the supervisor is not present to enforce it. Permits that are obtained retrospectively or not obtained at all are common enforcement findings.}
  \bitem{Acetylene cylinder backfire not managed}{Reverse flame entry into an acetylene cylinder during oxy-fuel cutting creates an unstable and potentially explosive cylinder. Operatives do not know the procedure for managing acetylene cylinder fires and use water on the cylinder, which is the wrong response.}
\end{itemize}

\example{Failure Pattern}{A heating contractor is installing a new gas main serving a residential development. The contractor engages a groundworker subcontractor to carry out the connection to the gas network main. The groundworker does not hold Gas Safe registration; the heating contractor has not checked this because the groundworker ``knows what he's doing around gas pipes''. The groundworker makes the main connection without pressure testing the joint. Gas escapes overnight into the ground, accumulates under the adjacent slab of a completed property, and migrates into the underfloor void. When the new resident's plumber opens the floor to connect an underfloor heating distribution manifold, a gas-air mixture is present; a small ignition source causes an explosion that seriously injures the plumber. The root cause investigation identifies unregistered gas work as the primary cause.}

%---------------------------- CHAPTER SUMMARY ----------------------------
\section{Chapter Summary}

\begin{table}[H]
\centering
\begin{tabularx}{\textwidth}{>{\raggedright\arraybackslash}p{3.2cm}X}
\toprule
\textbf{Assessment Element} & \textbf{Operational Requirement} \\
\midrule
Legal/Professional Basis & GSIUR 1998; DSEAR 2002; CDM 2015. All gas work must be conducted by a Gas Safe registered engineer; hot work permit system mandatory; one-hour post-work fire watch required. \\
Method & Apply the full five-step process and document inherent/residual risk. \\
Controls & Gas Safe registration verification; hot work permit with pre-work inspection; combustibles clearance; fire watch; LPG storage compliance; gas detection in enclosed spaces. \\
Cadence & Before each hot work permit; before each gas operative commences work; following any fire or gas escape; when new LPG equipment is introduced; minimum annually. \\
Assurance & Use audit, incident learning, and governance reporting to verify effectiveness. \\
\bottomrule
\end{tabularx}
\end{table}

\begin{itemize}
  \bitem{Key takeaway 1}{All gas work on construction sites must be carried out by a Gas Safe registered engineer with the correct category of registration for the work; verifying registration is the principal contractor's responsibility.}
  \bitem{Key takeaway 2}{Sub-surface fires in roof and floor build-ups are not detectable during the visible hot work period; a minimum one-hour post-work fire watch inspecting all adjacent areas above and below the work level is a non-negotiable control.}
  \bitem{Key takeaway 3}{LPG is heavier than air and accumulates in low-lying areas; LPG storage must comply with the Highly Flammable Liquids and LPG Regulations; LPG cylinders must never be stored below ground, in enclosed spaces, or in areas adjacent to heat sources.}
  \bitem{Key takeaway 4}{A hot work permit has value only if the pre-work inspection is genuinely conducted, combustibles are genuinely cleared, and the fire watch is genuinely maintained; a permit completed in retrospect or without the inspection and fire watch is worse than useless.}
\end{itemize}

\barquote{In Chapter~\ref{chap:Noise and Vibration Risk Assessment}, we turn from this assessment to the next domain in the Prudentia framework.}
